\section{Finite tests for the radical--CM compatibility defect}
\label{sec:lr-relative-gate}

Write $G=\widehat{\mathrm{GT}}$, $H=G_{\mathbb Q}$,
$k=\mathbb Q^{\mathrm{ab}}$, and $D=N_{\mathrm{rad}}$.
We use the established individual arithmetic quotients, the exact
radical--CM intersection, the signed CM-frame comparison, the normalized
signed-congruence identity, and the full hexagonal comparison of
\cite[Proposition~18]{H17} as inputs. Thus $G=DH$, and on $D$ the
multiplier and every rational Kummer cocycle vanish. The crossed defect is
\[
 d:G\longrightarrow C=\operatorname{Gal}(I/k),\qquad
 d(gh)=d(g)+\lambda_gd(h),\quad d|_H=0,
\]
with $C\simeq(\mathbb Z/12)^2\times\prod_{q\ne2,3}\mathbb Z/6$.
Nothing below assumes $d=0$, joint factorization, or $B/2B=0$.
Complete prime-two compatibility is a separate input when that
coordinate is removed from the final list of gates.

\subsection{An obstruction to a group-theoretic shortcut}

\begin{proposition}\label{prop:lr-abstract-countermodel}
Separate arithmetic quotient identifications do not imply vanishing
of a cyclotomically crossed defect. Even a trivial ordinary relative
abelian quotient does not force its order-three component to vanish.
\end{proposition}
\begin{proof}
Let $U=C_2$ act by negation on $V=\mathbb F_3$, and put
\[
 G_*=(V\oplus V)\rtimes U,\qquad
 H_*=\{(x,x,u):x\in V,\ u\in U\},\qquad A_*=V\rtimes U.
\]
The two projections $a_*,b_*:G_*\to A_*$ restrict to isomorphisms
on $H_*$. Hence $G_*=(\ker a_*)H_*=(\ker b_*)H_*$. Nevertheless
$d_*(x,y,u)=x-y$ is a nonzero crossed homomorphism, trivial on $H_*$.
Since $[G_*,G_*]=V\oplus V$, the abelianization of $G_*$ is $U$,
already covered by $H_*$. Its ordinary relative abelian quotient
is therefore zero. This is a countermodel to the proposed abstract
inference, not a construction of an element of GT.
\end{proof}

There is an exact cohomological reduction. Put
$G^+=\ker(\lambda\bmod3)$ and $H^+=H\cap G^+$, and choose complex
conjugation $c\in H$. Then $G=G^+\rtimes\langle c\rangle$.
Crossed homomorphisms to $\mathbb F_3(1)$ trivial on $H$ correspond
to homomorphisms
\[
 \chi:G^+\longrightarrow\mathbb F_3,\qquad
 \chi|_{H^+}=0,\qquad \chi(cgc^{-1})=-\chi(g).
\]
Restriction gives one direction; the inverse is $d(gc^e)=\chi(g)$.
Thus the relevant group is the minus eigenspace in the relative
abelianization of an index-two subgroup. Ordinary transfer applies
$1+c$ and kills this eigenspace; it does not show that it is zero.
Complex conjugation is not central in $G$, so the argument for a
central element acting nontrivially on coefficients cannot be used.
Odd conjugation weights all agree with the cyclotomic weight modulo
three and supply no weight mismatch.

By contrast, the two-primary image is the subgroup $d_2(D)$, of
exponent at most four, because $G=DH$. The action on
$d_2(D)/2d_2(D)$ is trivial, so $d_2$ induces ordinary relative
quadratic characters. A nonzero group of exponent four has a nonzero
quotient by twice itself. This observation does not prove that the
resulting characters vanish.

\subsection{The integral cubic-frame comparison}

Put $N=\ker(\lambda,\rho_2)$ and $\nu_g=B_g(\omega,\omega)$,
where $\omega^2+\omega+1=0$. If
$[f_g]=A_{f_g}(X,Y)[x,y]$ in $F'/F''$, then, for multiplier one,
\[
 B_g(X,Y)=1+(X-1)(Y-1)A_{f_g}(X,Y).
\]

\begin{proposition}\label{prop:lr-integral-cubic-beta}
For $g\in N$, an integral identification of the quadratic orders
identifies $T_g$ with either $(-1)^{a_2(g)}\nu_g$ or
$(-1)^{a_2(g)}\overline{\nu_g}$. For $g\in D$ the sign is one.
No uniform orientation is needed for cube membership or for the
scalar modulo-eight bit.
\end{proposition}
\begin{proof}
We specify the groups in the full comparison of
\cite[Proposition~18]{H17}. Put
\[
 \Gamma=B_3/\langle(uv)^3\rangle\simeq C_2*C_3,\qquad
 \Gamma_{1,1}=B_3/\langle(uv)^6\rangle,\qquad
 \Gamma_{1,2}=B_4/\langle(abc)^4\rangle,
\]
where $a,b,c$ are the standard chain generators of $B_4$.
The forgetting map
$\operatorname{forg}:\widehat\Gamma_{1,2}\to
\widehat\Gamma_{1,1}$ sends $(a,b,c)$ to $(u,v,u)$.
Its kernel is the completed punctured-torus group $\widehat\Pi$,
free on
\[
 \xi_1=ca^{-1},\qquad
 \xi_2=abac^{-1}b^{-1}a^{-1}.
\]
On $\widehat\Pi^{\mathrm{ab}}=\widehat{\mathbb Z}^2$,
$a,c$ act by $U$ and $b$ by $L$.

For $g=(1,f_g)$ let $\mathcal A_g$ be the braid action on
$\widehat\Gamma_{1,2}$:
\[
 \mathcal A_g(a)=a,\quad
 \mathcal A_g(b)=f_g(a^2,b^2)^{-1}b f_g(a^2,b^2),\quad
 \mathcal A_g(c)=P_g^{-1}cP_g,\quad
 P_g=f_g(ba^2b,c^2).
\]
Let $\widetilde A_g=A_g(a^2,b^2)$ be the evaluation of the
pure cubic conjugator $A_g\in\widehat F(x,y)$, and put
$\Psi_g=\operatorname{Ad}(\widetilde A_g^{-1})\mathcal A_g$.
Writing $J_3=\operatorname{forg}^{-1}(\langle uv\rangle)$,
the cited proposition gives an isomorphism
$\varphi:\widehat\Gamma\to J_3$, induced by a discrete
isomorphism and carrying $\overline{\Gamma'}$ to
$\widehat\Pi$, and an $h\in J_3$ such that
\[
 \Psi_g\varphi=\operatorname{Ad}(h)\varphi\Phi_g .
\]
Here $\overline{\Gamma'}$ is the closure of the commutator
subgroup of $\Gamma$, whose abelianization is the hexagonal
lattice $\widehat{\mathbb Z}[\omega]$.
We take this identity on the whole integral homology.

Let $M_g$ be the action of $\mathcal A_g$ on
$\widehat\Pi^{\mathrm{ab}}$. Equivariance of the forgetting map
and the established normalized signed-congruence identity imply
that $M_g$ commutes with $\mathrm{SL}_2(\widehat{\mathbb Z})$.
The centralizer of its two elementary unipotents consists of
scalars. Moreover $P_g\in\widehat\Pi$: after forgetting,
its two arguments are $(uv)^3u^{-2}$ and $u^2$, which commute,
and $f_g$ is a commutator word. Since
$ba^2b=(ab)^3a^{-2}$ and $c$ commutes with $a$, the element
$P_g$ commutes with $a$. Thus
$\mathcal A_g(\xi_1)=P_g^{-1}\xi_1P_g$, so $M_g$ fixes this
primitive homology class and $M_g=I$.

The corrected action on homology is consequently
$q(A_g)^{-1}=T_g$, since $D_g=I$ on $N$. The other side is
\[
 R_+^j C\nu_gC^{-1},\qquad j\in\mathbb Z/6,
\]
where $R_+=UL$ and the discrete homology identification $C$ sends the order-six
lattice rotation to $R_+$ or $R_+^{-1}$.
Lemma~15 of \cite{H17} gives
$\nu_g\equiv(-1)^{a_2(g)}\pmod4$.
Since $T_g\equiv I\pmod4$ and $R_+$ has exact order six modulo
four, $j=3a_2(g)\pmod6$. The two embeddings of the quadratic
order differ by its involution, proving the assertion.
On $D$, the established sign identity is
$a_2(g)=\rho_3^I(g)/2\bmod2=0$.

The identification of the lattice scalar with $B_g(\omega,\omega)$
is integral: $[x,y]$ maps to the translation $-3$ in the affine
triangle quotient, so $-3\nu_g=-3B_g(\omega,\omega)$.
Multiplication by three is injective in
$\widehat{\mathbb Z}[\omega]$, including at three, and can be cancelled.
\end{proof}

\subsection{The order-three gate}

\begin{lemma}\label{lem:lr-cubic-cm-quotient}
The maximal exponent-three subextension of $K_3^{\mathrm{ab}}/k$
is $k(q^{1/3}:q\text{ rational prime})$. The order-three part
of $d$ is therefore zero if and only if
$T_g\in N_3(4)^3$ for every $g\in D$.
\end{lemma}
\begin{proof}
Let $A=\operatorname{Gal}(K_3^{\mathrm{ab}}/K_3)$ and let $c$
act by complex conjugation. The full group over $\mathbb Q$ is
$A\rtimes\langle c\rangle$, whose closed derived subgroup is
$(c-1)A=\operatorname{Gal}(K_3^{\mathrm{ab}}/k)$.
On the pro-three part the idempotents $(1\pm c)/2$ split $A$,
and $(c-1)A$ is precisely its minus part.
Every cubic character of this subgroup thus extends to $A$,
with $c$ acting by negation on the cyclic cubic quotient.

Kummer theory writes the corresponding field as
$K_3(\alpha^{1/3})$, where the class of $\alpha$ in
$K_3^\times/(K_3^\times)^3$ is fixed by conjugation: inversion
on the Galois quotient is already accounted for by the action
on $\mu_3$. A fixed class comes from $\mathbb Q$ because
$[N_{K_3/\mathbb Q}(\alpha)]=[\alpha]^2$, and squaring again
recovers $[\alpha]$ in this exponent-three group.
Conversely rational cube roots give such extensions.
Character separation proves the simultaneous field statement.
The normalized cubic torus is an isomorphic model of the Galois
group over $k$; its quadratic sign correction has no effect on
the quotient by cubes. On $D$ the radical restriction is trivial.
\end{proof}

\begin{lemma}\label{lem:lr-wild-three-cube}
For every $g\in D$, the three-adic component of $\nu_g$ is a cube
in the norm-one units of $\mathbb Z_3[\omega]$.
\end{lemma}
\begin{proof}
Put $\pi=\omega-1$. The class-two three-cycle relation gives
$3A_{f_g}(1,1)=0$, hence $A_{f_g}(1,1)=0$.
Symmetry gives a three-adic expansion
\[
 A_{f_g}(1+S,1+T)=a(S+T)+O((S,T)^2).
\]
In the degree-three cyclic-cover homology modulo three one has
$Y=1$, $X^3=1$, and $[x,y]\mapsto(X-1)y_0$.
Since $(X-1)^3=0$ in $\mathbb F_3[C_3]$, the image of $f_g$ is
$a(X-1)^2y_0=a(1+X+X^2)y_0$.
Its $y_0$ coefficient is $a$, proving
$\rho_3^I(g)\equiv-a\pmod3$. This is zero on $D$. Thus
\[
 \nu_g=1+\pi^2A_{f_g}(\omega,\omega)
      =1+2a\pi^3+O(\pi^4)\equiv1\pmod{\pi^4}.
\]
The unit has norm one by
Proposition~\ref{prop:lr-integral-cubic-beta}, and it is one
modulo three by \cite[Lemma~15]{H17}.
The logarithm identifies the norm-one units in $1+\pi^3O_3$
with $3\sqrt{-3}\,\mathbb Z_3$. Congruence modulo $\pi^4$
forces the logarithm into $9\sqrt{-3}\,\mathbb Z_3$, whose
valuation is five. This is three times the previous lattice;
exponentiation proves cube membership. The possible $\mu_3$
factor was already excluded by congruence to one modulo three.
\end{proof}

\begin{theorem}\label{thm:lr-tame-cubic-gate}
The order-three component of $d$ vanishes if and only if, for
all $g\in D$ and primes $\ell\ge5$,
\begin{equation}
 B_{g,\ell}(\omega,\omega)^{(\ell-\chi_{-3}(\ell))/3}=1
 \quad\text{in the residue norm-one torus}.
 \label{eq:lr-tame-cubic-test}
\end{equation}
Use a primitive cube root in $\mathbb F_\ell$ in the split
case and in $\mathbb F_{\ell^2}$ in the inert case.
\end{theorem}
\begin{proof}
Cubing is bijective on the pro-two factor $N_{3,2}(4)$.
The three-component is handled by the preceding lemma.
For $\ell\ge5$, local reduction has a pro-$\ell$ kernel and a
cyclic norm-one image of order $\ell-\chi_{-3}(\ell)$.
Cubing is bijective on the kernel, and its image on the residue
group is exactly the kernel of the displayed power.
Apply Proposition~\ref{prop:lr-integral-cubic-beta} and
Lemma~\ref{lem:lr-cubic-cm-quotient}. Inversion preserves the tests.
For the product, the cube image is the product of the local
cube images, so no simultaneous condition has been discarded.
\end{proof}

The Anderson relation $(A_\ell)$ is sufficient for
\eqref{eq:lr-tame-cubic-test}. In its beta form
\cite[Remark~2.4]{Furusho2003}, specialize at $s=1/3$ and
reduce in characteristic $\ell$. The $\ell$th roots in the
denominator all reduce to one, where $B_g(1,\omega)=1$.
The numerator contains $B_g(\omega,\omega)$ exactly
$(\ell-\chi_{-3}(\ell))/3$ times; all other factors are boundary
values equal to one. The arithmetic right side is one on $D$.
This implication does not assert $(A_\ell)$ for all of GT.

\subsection{A formal obstruction to removing the tame tests}

\begin{proposition}\label{prop:lr-formal-beta-obstruction}
Beta associativity, unitarity, all dyadic Anderson identities,
vanishing integral Kummer coordinates, and the logarithmic
derivatives of the Anderson identities do not, as formal integral
Alexander identities, imply all tests
\eqref{eq:lr-tame-cubic-test}.
\end{proposition}
\begin{proof}
Fix $\ell\ge5$. By the full affine radical image choose
$\sigma\in H^{\mathrm{cyc}}$ whose sole nonzero radical
translation is $\rho_\ell^I(\sigma)=2e_3$, with
$e_3\in\widehat{\mathbb Z}$ the coefficient-prime-three
idempotent. In particular $\rho_2^I(\sigma)=\rho_3^I(\sigma)=0$
and every $\ell$-adic Kummer coordinate is zero.
The arithmetic $(A_\ell)$ relation at $s=1/3$ gives
\[
 B_{\sigma,\ell}(\omega,\omega)^{
          (\ell-\chi_{-3}(\ell))/3}=\omega^{-2\ell}\ne1
 \quad\bmod\ell.
\]
Take the coefficient-$\ell$ component of $A_{f_\sigma}$ and
set every other coefficient-prime component to zero.
Its beta unit has component $B_{\sigma,\ell}$ at $\ell$ and
one elsewhere. This is integral in the completed group algebra.
Associativity, unitarity, and the dyadic identities hold
componentwise, because they hold for $\sigma$ and the identity.
Every cyclic-cover coefficient is a
$\widehat{\mathbb Z}$-linear functional on $F'/F''$.
Thus every Kummer coordinate of the patched data is zero.
The logarithmic derivatives here are taken in the coefficient-$p$
formal variable near each torsion specialization, for every
coefficient prime $p$. At $p=\ell$, write the variable as
$\zeta(1+S)$, with $\zeta$ of order prime to $\ell$.
The arithmetic Anderson discrepancy for $\sigma$ has principal
factor $(1+S)^{n\rho_{n,\ell}^I(\sigma)}=1$; its remaining
factor is constant torsion on that specialization.
Its logarithmic derivative is therefore zero. At every other
coefficient prime the patched unit is one. This proves all the
stated derivative identities while leaving the displayed
nontrivial constant $\omega^{-2\ell}$ intact.
The displayed tame test still fails.

This constructs formal Alexander data, not a lift satisfying the
full profinite pentagon. It also does not assert the full
nonabelian $(K_n)$ identities for a lift; these contain information
beyond the numerical Kummer and logarithmic-derivative consequences.
The example isolates a missing constraint in the proposed formal
inference without supplying a counterexample inside GT.
\end{proof}

\subsection{The quartic coordinates}

\begin{lemma}\label{lem:lr-quartic-cm-quotient}
The maximal elementary-two subextension of
$K_4^{\mathrm{ab}}/k$ is
\[k(q^{1/4}:q\text{ rational prime}).\]
\end{lemma}
\begin{proof}
Use CM reciprocity in the primary-unit realization
\[
 \operatorname{Gal}(K_4^{\mathrm{ab}}/K_4)
 \simeq\{u\in\widehat O_4^\times:
                    u_2\equiv1\pmod{(1+i)^3}\},
\]
whose commutator image is $N_4(4)$.
The local two-component is isomorphic by logarithm to
$4i\mathbb Z_2$ and has one quadratic character.
Each odd component has one quadratic character, since its
residue norm-one group is cyclic of order
$q-\chi_{-4}(q)$ and its principal kernel is pro-$q$.

The character of $2^{1/4}$ is supported only at two and is
nontrivial over $k$, so it spans the first character space.
For odd $q$, $K_4(q^{1/4})/K_4$ has conductor supported at
two and $q$, and its tame character at $q$ has order four.
On norm-one units it restricts to the unique nontrivial
quadratic character: writing such a unit as $u/\bar u$,
complex conjugation inverts the Kummer character, so its value
is the square of the order-four tame character of $u$.
The local ratio map is surjective.
Thus the restriction is the tame $q$-character plus a possible
two-component character. Conductor support and the primary-unit
form of reciprocity exclude every other local component.
These characters, together with the character of $2^{1/4}$,
span the continuous quadratic character space of the product.
Character separation proves the result. The spanning argument
uses local support, not a comparison of infinite dimensions.
\end{proof}

On $N$ one has $W_{g,2}\equiv I\pmod8$. Indeed
$\beta_g\in F'$ and $D_g=I$, and modulo eight the commutator
subgroup generated by $U^2,L^2$ is $\{I,5I\}$.
The reduction $5I$ cannot lift to a norm-one element $a+bi$:
if $a\equiv5$ and $b\equiv0\pmod8$, then
$a^2+b^2\equiv9\pmod{16}$. This observation is subsumed by
the separate stronger prime-two compatibility theorem.
Consequently the local two-component of $W_g$ is a square.
At three, the established sign identity on $D$ places
$W_{g,3}\bmod3$ in $\{I,-I\}$, the square subgroup of the
cyclic order-four residue torus.
Lemma~\ref{lem:lr-quartic-cm-quotient} now gives the exact
simultaneous reduction: all rational quartic compatibilities
hold if and only if
\begin{equation}
 W_{g,\ell}^{(\ell-\chi_{-4}(\ell))/2}=1
 \quad\text{in the residue norm-one torus}
 \qquad(g\in D,\ \ell\ge5).
 \label{eq:lr-tame-quartic-test}
\end{equation}
Proposition~\ref{prop:tame-harmonic-word-gate} expresses these
tests directly as finite word evaluations by the normalized
harmonic identity. Their vanishing remains open.

\subsection{The exceptional eighth root of three}

Put $E_3:y^2=x^3+1$, $\alpha=3^{1/4}$, and
$v=(2+\sqrt3)^{1/4}$, using positive roots.

\begin{lemma}\label{lem:lr-cubic-eight-field}
One has $k(E_3[8])=k(\alpha,v)$.
The scalar automorphism $[-1]$ of this division field fixes $v$.
On $H^{\mathrm{cyc}}$, if
\[
 T_{\sigma,2}\equiv(1+4\tau(\sigma))I\pmod8,
 \qquad \tau(\sigma)\in\mathbb F_2,
\]
then $\sigma(v)/v=(-1)^{\tau(\sigma)}$.
\end{lemma}
\begin{proof}
Write $e_j=-\omega^j$, for $j=0,1,2$.
The abscissae of the non-two-torsion points in $E_3[4]$ are
$\omega^j(-1+\sqrt3)$ and $\omega^j(-1-\sqrt3)$.
The established four-division calculation is
$k(E_3[4])=k(\alpha)$.
To halve $P=(x_P,y_P)$ one adjoins square roots of $x_P-e_j$
with product of the prescribed sign $-y_P$.
Conversely these square roots belong to the field of a half
$Q=(x_Q,y_Q)$, because duplication gives
\[
 x_P-e_j=
 \left(\frac{(x_Q-e_j)^2-(e_j-e_{j+1})(e_j-e_{j+2})}
 {2y_Q}\right)^2.
\]
Indices are modulo three. This supplies both directions of
the square-root criterion inside the full division field.

For $x_P=-1+\sqrt3$, set $A=\sqrt3-1+\omega$.
Since $v^{-2}=\sqrt{2-\sqrt3}$, direct multiplication gives
\[
 A=\zeta_{24}^{\,5}\sqrt3\,v^{-2},\qquad
 \sqrt A=\zeta_{48}^{\,5}\alpha/v
 \quad\text{up to sign}.
\]
Hence $v\in k(E_3[8])$.
Conversely every square root needed to halve the displayed
four-torsion abscissae is a root of unity times
$\alpha$, $\alpha/v$, or $\alpha v$.
For the minus choice the relevant identity is
\[
 -\sqrt3-1+\omega=\zeta_{24}^{\,11}\sqrt3\,v^2.
\]
Conjugating and multiplying by roots of unity gives all cases.
Thus the halving formulas place all of $E_3[8]$ in $k(\alpha,v)$.

The two quadratic extensions $k(\alpha)$ and $k(v)$ are
nontrivial and distinct. The positive numbers $\alpha,v,v/\alpha$
each have a nonreal conjugate, whereas any real element of
$k$ generates a totally real abelian field.
For $v/\alpha$, the automorphism of
$\mathbb Q(i,\sqrt2,\sqrt3)$ fixing $\sqrt3$ and changing
the sign of $\sqrt2$ negates $(v/\alpha)^2$.
The same criterion, or the defining polynomials, handles
$\alpha$ and $v$. Hence $[k(\alpha,v):k]=4$.

Under $[-1]$, the four-torsion ordinate changes sign,
so $\alpha$ changes sign. The recovered square root in
the duplication formula also changes sign when $Q$ is
replaced by $-Q$. Their quotient $v$ is therefore fixed.
The unnormalized cubic norm-one image modulo eight,
one modulo two, consists of the scalar units $1,3,5,7$.
The scalar $-1=7$ fixes $v$, and the other nontrivial
automorphism fixing $E_3[4]$, the scalar $5$, changes its sign.
The normalization from the Tate representation to $T$
multiplies by a scalar sign, which fixes $v$.
Its image modulo eight is $\{I,5I\}$, proving the assertion.
\end{proof}

For $g\in D$ define $w_3(g),\tau_3(g)\in\mathbb F_2$ by
\[
 W_{g,3}\bmod3=(-I)^{w_3(g)},\qquad
 T_{g,2}\bmod8=(1+4\tau_3(g))I.
\]
The first uses the known vanishing of the individual frame
parity; the second follows from norm one and congruence one
modulo four. Both are homomorphisms on $D$.

\begin{proposition}\label{prop:lr-exceptional-three-character}
The character $\chi_3=w_3+\tau_3$ on $D$ is $G$-invariant
and zero on $D\cap H$. It extends uniquely to an ordinary
continuous quadratic character on $G$ that is zero on $H$.
Its vanishing is exactly the remaining compatibility for
$3^{1/8}$, once the known quartic coordinate of three is imposed.
Moreover
\[
 \tau_3(g)=\frac{B_{g,2}(\omega,\omega)-1}{4}\pmod2,
\]
where the right side is scalar, in
$\mathbb F_2\subset\mathbb F_4$.
\end{proposition}
\begin{proof}
The harmonic frame extends to a matrix representation on the
whole gentle group, so conjugation fixes its scalar classes
$\{I,-I\}$ modulo three. The full cubic-frame conjugation
formula fixes the scalar classes $\{I,5I\}$ modulo eight.
Thus both bits are $G$-invariant.

Use the arithmetic witness in the intersection theorem:
an ordinate $y$ of a three-torsion point of
$E_4:y^2=x^3-x$ satisfies
\[
 y^4=\frac49\sqrt3(2+\sqrt3),\qquad
 z=\frac{y}{\sqrt{2/3}\,v},\qquad z^8=3.
\]
For $\sigma\in D\cap H$, the harmonic residue frame sends
$y$ to $(-1)^{w_3(\sigma)}y$, while
Lemma~\ref{lem:lr-cubic-eight-field} gives the sign
$(-1)^{\tau_3(\sigma)}$ on $v$.
Since $\sigma$ fixes every rational radical, it fixes $z$,
so the sum of the two bits is zero.

The extension is $\chi_3(dh)=\chi_3(d)$ for $d\in D,h\in H$.
Invariance and vanishing on the intersection prove
well-definedness and multiplicativity. Continuity follows
from the compact surjection $D\times H\to G$.
On $D$, the radical coordinate and the quartic coordinate of
three are zero, and the sign on the eighth root is exactly
this sum, proving its interpretation as the remaining gate.

Finally the integral cubic-frame comparison identifies $T_g$
with $B_g(\omega,\omega)$ or its conjugate.
Modulo eight a norm-one unit that is one modulo four is
the scalar $1$ or $5$. Conjugation changes neither, proving
the beta formula. Its vanishing on all of $G$ is proved in
Section~\ref{sec:normalized-cubic-identity}.
\end{proof}

\subsection{Status}

After the separate prime-two compatibility theorem, joint
compatibility is reduced to the tame cubic tests
\eqref{eq:lr-tame-cubic-test}, the tame harmonic quadratic
tests \eqref{eq:lr-tame-quartic-test}, and the single relative
character $\chi_3$. The wild three-adic cubic part is eliminated.
The two subsequent frame comparisons also give $\chi_3=0$,
so only the tame tests remain.
The exceptional eighth root of three is represented by a
concrete finite combination of cubic eight-division and
harmonic three-division frames. The individual quotients,
transfer, and the formal beta identities considered here do
not force the remaining tests to vanish. Full joint
factorization remains unproved.
